\section{Considerações finais}
\label{sec:finais}

No estágio atual do trabalho, a resposta à pergunta de pesquisa é parcial. A dissertação organizou um protocolo que articula recuperação documental, geração de respostas, avaliação humana e julgamento automatizado, e o aplicou preliminarmente a um sistema RAG voltado aos discursos da 56\textordfeminine{} Legislatura do Senado Federal. Como o experimento principal ainda não foi realizado, as considerações a seguir se restringem às contribuições documentadas e às etapas necessárias para avaliar o protocolo em condições controladas.

\subsection{Síntese das contribuições até o momento}

A pesquisa sistematizou técnicas de avaliação de RAG e as organizou em um protocolo que mantém distintas, mas articuladas, a recuperação, a geração, a verificabilidade das fontes, a avaliação humana, o julgamento automatizado e a análise qualitativa de erros. O protocolo especifica métricas de recuperação, uma rubrica de cinco dimensões e regras de interpretação para H1--H3. Essa contribuição é metodológica e permanece como proposta a ser avaliada na etapa principal; sua especificação não equivale à validação empírica do protocolo.

O estudo também caracterizou o corpus documental e documentou a preparação de uma bateria de 200 perguntas distribuídas em 16 categorias, da qual foi selecionado um piloto estratificado de 32 itens. Na execução de viabilidade reportada, os 32 itens foram processados sem erro, com média automatizada de 9,094 em escala de 0 a 10. O registro sinalizou risco alto de mistura de autores em 28 itens. Esses achados demonstram que o fluxo pôde ser executado e produziu material para inspeção; não medem a qualidade geral das respostas nem o desempenho da recuperação.

Assim, a resposta provisória à pergunta de pesquisa é que as técnicas podem ser organizadas em um protocolo aplicado ao caso empírico, desde que suas dimensões sejam avaliadas separadamente e relacionadas por evidências rastreáveis. A execução disponível demonstra a viabilidade operacional de uma aplicação-piloto, mas ainda não permite concluir sobre a eficácia comparativa das técnicas ou sobre as hipóteses da dissertação.

\subsection{Limitações identificadas}

A principal limitação é a ausência do experimento principal: ainda não há conjunto de referência validado, comparação controlada entre BM25, recuperação vetorial e híbrida, nem resultados de avaliação humana independente. Por isso, não é possível calcular as métricas de recuperação previstas, comparar respostas por categoria ou estimar concordância entre avaliadores humanos e juiz automatizado. H1, H2 e H3 permanecem sem avaliação empírica conclusiva.

Os resultados de geração disponíveis provêm de uma rubrica exploratória e de uma execução em que \texttt{gpt-5-nano} atuou como gerador e juiz. A execução registrada também difere em dois enunciados da versão da bateria atualmente preservada no projeto. Além disso, temperatura, \textit{top-p} e limite de tokens não foram informados, o que limita a reprodução exata das condições. A bateria do piloto contém apenas dois itens por categoria e não representa uma amostra probabilística de consultas reais. Por fim, o estudo se restringe ao corpus de discursos da 56\textordfeminine{} Legislatura e às configurações documentadas; seus resultados não podem ser generalizados para outros acervos ou sistemas RAG.

\subsection{Próximos passos da pesquisa}

As etapas seguintes devem ocorrer na seguinte ordem:

\begin{enumerate}
  \item Congelar a bateria e sua versão analítica, resolver as divergências de enunciado, confirmar a respondibilidade e validar, para cada pergunta, as evidências documentais suficientes, inclusive nos casos parcialmente respondíveis e sem suporte.
  \item Fixar a configuração do sistema e da segmentação e registrar modelos, \textit{embeddings}, parâmetros de recuperação e geração, \textit{prompts}, versões e sementes aplicáveis.
  \item Executar BM25, recuperação vetorial e híbrida sob as mesmas perguntas e evidências de referência; calcular as métricas compatíveis com a granularidade dos julgamentos e analisar separadamente documentos, trechos e afirmações.
  \item Aplicar a rubrica principal às respostas, realizar avaliação humana independente e comparar seus julgamentos com os do avaliador automatizado, preservando notas individuais, divergências e casos críticos.
  \item Aplicar as regras de decisão de H1--H3 aos resultados, discutir as condições e os limites observados e avaliar o protocolo como artefato metodológico, registrando eventuais revisões decorrentes de sua aplicação.
\end{enumerate}

Somente após essas etapas será possível substituir a resposta provisória desta versão por conclusões empíricas sobre a complementaridade e os limites das técnicas avaliativas.

\subsection{Cronograma de execução}

\label{sec:cronograma}

A Tabela~\ref{tab:cronograma} reproduz as etapas, entregas ao orientador e resultados esperados do cronograma registrado.

\begin{landscape}
\begingroup
\small
\setlength{\tabcolsep}{4pt}
\renewcommand{\arraystretch}{1.15}
\begin{longtable}{p{0.12\linewidth}p{0.28\linewidth}p{0.30\linewidth}p{0.24\linewidth}}
\caption{Cronograma de atividades para dissertação}\label{tab:cronograma}\\
\toprule
\textbf{Período} & \textbf{Atividade} & \textbf{Entrega para o orientador} & \textbf{Resultado esperado} \\
\midrule
\endfirsthead
\multicolumn{4}{c}{\tablename~\thetable{} -- continuação}\\
\toprule
\textbf{Período} & \textbf{Atividade} & \textbf{Entrega para o orientador} & \textbf{Resultado esperado} \\
\midrule
\endhead
\midrule
\multicolumn{4}{r}{Continua na página seguinte}\\
\endfoot
\bottomrule
\endlastfoot
ago/2026 & Alinhar expectativas com o orientador e validar pergunta de pesquisa, objetivos, hipóteses, escopo e coleta de dados & Pré-projeto revisado; pauta de decisões metodológicas & Escopo da qualificação confirmado \\
set/2026 & Consolidar o referencial teórico e a metodologia para qualificação & Capítulos de introdução, referencial teórico, hipóteses e metodologia em versão expandida & Estrutura da qualificação validada \\
out/2026 & Fechar o texto para qualificação e preparar a apresentação & Versão quase final da qualificação; slides de apresentação de até 15 minutos & Aprovação do orientador para envio à banca \\
Até 07/11/2026 & Enviar formalmente o material à banca de qualificação & Documento de qualificação enviado à banca & Cumprimento da antecedência mínima de 10 dias \\
17/11/2026 & Realizar a qualificação da dissertação & Apresentação oral e defesa do projeto perante a banca & Recebimento das recomendações da banca \\
18/11 a 15/12/2026 & Organizar e incorporar inicialmente as recomendações da banca & Matriz pós-qualificação com ajustes obrigatórios, recomendados e opcionais & Plano de ajustes pós-qualificação validado \\
16/12/2026 a 31/01/2027 & Preparar a execução empírica durante a licença & Corpus localizado; ambiente técnico preparado; checklist de execução; método ajustado após a banca & Condições técnicas e metodológicas preparadas para início da licença \\
01/02 a 14/02/2027 & Planejar o início do afastamento e organizar os instrumentos de acompanhamento & Plano de trabalho do primeiro período da licença; agenda de reuniões; repositório organizado & Prioridades da licença definidas com o orientador \\
15/02 a 21/02/2027 & Aprofundar a revisão bibliográfica e documental dirigida e realizar ajustes finais pós-qualificação & Matriz teórica revisada; lista final de referências essenciais & Aprofundamento da revisão bibliográfica e documental \\
22/02 a 28/02/2027 & Consolidar o referencial teórico e o percurso metodológico & Capítulos de introdução, referencial, hipóteses e metodologia revisados & Consolidação do referencial teórico e metodológico \\
01/03 a 07/03/2027 & Congelar o corpus e preparar tecnicamente o pipeline & Manifesto de dados; regras de inclusão e exclusão; parâmetros preliminares do pipeline RAG & Percurso metodológico consolidado e corpus versionado \\
08/03 a 14/03/2027 & Construir a bateria de perguntas e o conjunto de referência & Bateria de perguntas por categoria; evidências documentais esperadas & Projeto analítico consolidado \\
15/03 a 19/03/2027 & Validar o material produzido no primeiro período da licença & Sumário estruturado da dissertação; capítulos iniciais revisados; protocolo avaliativo v1 & Versão preliminar dos capítulos iniciais e da metodologia \\
20/03 a 31/03/2027 & Revisar, com o orientador, o material produzido no primeiro período da licença & Devolutiva sobre corpus, perguntas, protocolo e capítulos & Ajustes críticos identificados antes da execução empírica \\
01/04 a 04/04/2027 & Preparar a execução empírica final & Lista fechada de ajustes antes da execução do protocolo & Protocolo pronto para aplicação \\
05/04 a 11/04/2027 & Executar o pipeline RAG e coletar os resultados & Resultados brutos: evidências recuperadas e respostas geradas & Execução de testes, experimentos e análises preliminares \\
12/04 a 18/04/2027 & Realizar a avaliação humana e a avaliação baseada no paradigma LLM-as-a-judge & Planilha de julgamentos humanos; resultados automatizados; registro dos parâmetros de execução & Sistematização dos resultados parciais \\
19/04 a 23/04/2027 & Analisar os resultados e redigir os capítulos empíricos & Capítulo de resultados v1; análise preliminar das hipóteses H1, H2 e H3 & Minuta dos capítulos de desenvolvimento e análise inicial dos resultados \\
24/04 a 26/04/2027 & Revisar a versão final para a banca de defesa & Versão completa da dissertação para revisão conclusiva do orientador & Autorização para depósito da versão de defesa \\
Até 27/04/2027 & Depositar ou enviar a versão à banca de defesa & Dissertação enviada à banca & Cumprimento da antecedência mínima de 10 dias para a defesa \\
28/04 a 30/04/2027 & Preparar a apresentação oral da defesa & Slides; roteiro de apresentação de 20 minutos; simulação de arguição & Defesa preparada \\
07/05/2027 & Defender a dissertação & Apresentação oral, arguição e deliberação da banca & Aprovação, aprovação com ajustes ou reformulação \\
08/05 a 15/05/2027 & Organizar as recomendações da banca & Matriz de correções pós-defesa & Plano de ajustes pós-defesa definido \\
16/05 a 30/05/2027 & Incorporar os ajustes iniciais recomendados pela banca & Versão pós-defesa v1 para o orientador & Material preparado para consolidação final durante o terceiro período da licença \\
31/05 a 06/06/2027 & Incorporar as recomendações da banca aos capítulos afetados & Versão revisada dos capítulos afetados & Análise final dos resultados e discussão incorporadas \\
07/06 a 13/06/2027 & Consolidar a conclusão, as recomendações e os limites da pesquisa & Conclusão revisada; recomendações de uso cauteloso; limites da pesquisa & Capítulos finais consolidados \\
14/06 a 20/06/2027 & Revisar o texto do ponto de vista acadêmico, metodológico e bibliográfico & Versão completa revisada & Texto revisado do ponto de vista acadêmico, metodológico e bibliográfico \\
21/06 a 25/06/2027 & Normalizar a dissertação conforme normas acadêmicas & Versão normalizada da dissertação & Adequação às normas acadêmicas e formatação final \\
26/06 a 30/06/2027 & Fechar a versão final e preparar o depósito institucional & Versão final depositável; plano do artigo derivado & Dissertação concluída em versão final, texto revisado e normalizado, versão pronta para depósito institucional \\
\end{longtable}
\noindent Fonte: cronograma de trabalho do pesquisador, em \path{cronograma/cronograma.tex}; datas planejadas, sujeitas a confirmação.
\endgroup
\end{landscape}
